\documentclass[11pt]{article}
\input{../../shared/project_style.tex}
\rhead{Basics 1 Textbook Draft}

\begin{document}
\DocTitle{Basics 1: Units, Dimensions, and Normalization}{I - Mathematical and computational foundations}{A practical language for checking equations and constructing dimensionless plasma models}

\begin{overviewbox}
\textbf{Textbook draft, version 1.0.} Units specify a measurement convention, dimensions specify the physical type of a quantity, and normalization exposes the dimensionless parameters that organize a model. This chapter is designed for a reader with no prior plasma-physics training. It distinguishes exact identities from physical approximations and connects the central equations to later simulation modules.
\end{overviewbox}

\ProjectTOC

\section{Learning objectives}
After completing this note, the reader should be able to:
\begin{itemize}[leftmargin=1.6em]
\item Distinguish dimensions from units and convert among common plasma-physics units.
\item Check dimensional homogeneity of algebraic, differential, and integral equations.
\item Choose characteristic scales for length, time, velocity, density, pressure, magnetic field, and frequency.
\item Nondimensionalize conservation laws and identify the resulting control parameters.
\item Explain how normalization affects numerical conditioning and physics-informed loss functions.
\item Document units and normalization at every software interface.
\end{itemize}

\section{Prerequisites and reading strategy}
\begin{itemize}[leftmargin=1.6em]
\item Elementary algebra, scientific notation, and differentiation.
\end{itemize}
On a first reading, emphasize physical meaning and dimensional checks. On a second reading, reproduce the derivations. Attempt the computational laboratory only after the limiting cases are understood.

\section{Physical quantities, base dimensions, and SI units}

A physical quantity is a numerical value multiplied by a unit. The same magnetic field may be written as
\begin{equation}
B=2.5\ \si{T}=2.5\times10^{4}\ \si{G},
\end{equation}
but its physical dimension is unchanged. In SI, the base dimensions needed here are mass $M$, length $L$, time $T$, electric current $I$, and thermodynamic temperature $\Theta$. We write
\begin{equation}
[Q]=M^aL^bT^cI^d\Theta^e
\end{equation}
to classify a quantity $Q$. Dimensions are independent of the human choice of meter, centimeter, second, or microsecond.

Examples are
\begin{align}
[\rho]&=ML^{-3},&
[p]&=ML^{-1}T^{-2},\\
[\mathbf u]&=LT^{-1},&
[\mathbf B]&=MT^{-2}I^{-1},\\
[\mathbf E]&=MLT^{-3}I^{-1},&
[\mu_0]&=MLT^{-2}I^{-2}.
\end{align}
The magnetic-energy density $B^2/(2\mu_0)$ has dimension $ML^{-1}T^{-2}$, which is exactly the dimension of pressure. That equality is the dimensional reason magnetic pressure can enter the MHD momentum equation beside gas pressure.


\section{Dimensional homogeneity as an error detector}

Every term that is added to another term must have the same dimension. Consider
\begin{equation}
\rho\left(\pdv{\mathbf u}{t}+\mathbf u\cdot\nabla\mathbf u\right)
=-\nabla p+\frac{1}{\mu_0}(\nabla\times\mathbf B)\times\mathbf B.
\end{equation}
The inertial term scales as $\rho U^2/L$, the pressure force as $p/L$, and the magnetic force as $B^2/(\mu_0L)$. All are force per volume. If a derivation accidentally uses $B/\mu_0L$ rather than $B^2/\mu_0L$, dimensional analysis rejects it immediately.

Dimensional consistency is necessary but not sufficient. The expression $mv^2$ and $2mv^2$ have the same dimension, but only one coefficient may be correct. Dimensional analysis also cannot detect a wrong sign or an omitted dimensionless factor such as $2\pi$. It should therefore be the first verification step, not the last.


\section{Characteristic scales}

A characteristic scale is a representative magnitude chosen to organize a problem. Introduce
\begin{equation}
\mathbf x=L_0\hat{\mathbf x},\quad
t=t_0\hat t,\quad
\mathbf u=U_0\hat{\mathbf u},\quad
\rho=\rho_0\hat\rho,
\end{equation}
and similarly
\begin{equation}
\mathbf B=B_0\hat{\mathbf B},\quad
p=p_0\hat p,\quad
\omega=\omega_0\hat\omega.
\end{equation}
For magnetized-fluid dynamics, natural choices include
\begin{equation}
V_{A0}=\frac{B_0}{\sqrt{\mu_0\rho_0}},
\qquad
t_A=\frac{L_0}{V_{A0}},
\qquad
p_{B0}=\frac{B_0^2}{\mu_0}.
\end{equation}
An orbit code may instead use a gyroperiod and gyroradius. A transport code may use a confinement time and minor radius. There is no universally correct normalization; there is only a normalization that is explicit, internally consistent, and appropriate to the dominant scales.


\section{Nondimensionalizing representative equations}

The continuity equation
\begin{equation}
\pdv{\rho}{t}+\nabla\cdot(\rho\mathbf u)=0
\end{equation}
becomes
\begin{equation}
\frac{\rho_0}{t_0}\pdv{\hat\rho}{\hat t}
+\frac{\rho_0U_0}{L_0}\hat\nabla\cdot(\hat\rho\hat{\mathbf u})=0.
\end{equation}
Choosing $t_0=L_0/U_0$ removes the prefactors. For magnetic induction with magnetic diffusivity $\eta_m$,
\begin{equation}
\pdv{\mathbf B}{t}
=\nabla\times(\mathbf u\times\mathbf B)+\eta_m\nabla^2\mathbf B,
\end{equation}
the normalized equation is
\begin{equation}
\pdv{\hat{\mathbf B}}{\hat t}
=\hat\nabla\times(\hat{\mathbf u}\times\hat{\mathbf B})
+\frac{1}{R_m}\hat\nabla^2\hat{\mathbf B},
\qquad
R_m=\frac{U_0L_0}{\eta_m}.
\end{equation}
The magnetic Reynolds number compares field advection with magnetic diffusion. In the momentum equation, the coefficient
\begin{equation}
\frac{B_0^2}{\mu_0\rho_0U_0^2}=\frac{1}{M_A^2}
\end{equation}
introduces the Alfv\'en Mach number $M_A=U_0/V_A$.


\section{Frequency, density, and temperature conventions}

Angular frequency and cyclic frequency satisfy
\begin{equation}
\omega=2\pi f.
\end{equation}
A value in radians per second is not numerically equal to the value in hertz. Number density and mass density are related by
\begin{equation}
\rho=\sum_s m_sn_s.
\end{equation}
For an Alfv\'en speed, using electron density alone would omit the dominant ion inertia.

The electron-volt is an energy unit:
\begin{equation}
1\ \si{eV}=1.602176634\times10^{-19}\ \si{J}.
\end{equation}
Plasma temperature is often reported as $k_BT$ in electron-volts. A field named ``temperature\_eV'' normally stores an energy per particle, not kelvin. The code and the notes must state the convention.


\section{Normalization in numerical simulation and PINNs}

Floating-point calculations are more robust when important variables are of comparable magnitude. Normalized variables also make residual tolerances meaningful across modules.

In a PINN, suppose
\begin{equation}
\mathcal L=\mathcal L_{\rm mom}+\mathcal L_{\nabla\cdot B}.
\end{equation}
If dimensional residuals differ by many orders of magnitude, optimization may effectively ignore one equation. Dimensionless residuals make the relative weights interpretable, although adaptive weighting may still be needed.

\begin{physicsbox}
A reproducible data product states the physical unit, normalization scale, sign convention, and whether the stored value is dimensional or normalized.
\end{physicsbox}


\section{Computational laboratory}

Create matching Python and MATLAB scripts that read $B_0$, ion number density, ion mass, and major radius from JSON. Compute mass density, Alfv\'en speed, Alfv\'en time, angular frequency, and cyclic frequency. Double $B_0$ and quadruple $\rho_0$ and verify the expected scalings. Add tests that reject negative density and ambiguous frequency metadata.


\section{Summary}
\begin{itemize}[leftmargin=1.6em]
\item Dimensions classify physical quantities independently of the chosen unit system.
\item Dimensional homogeneity is a powerful necessary check on every equation.
\item Characteristic scales expose the dimensionless parameters controlling a model.
\item Frequency, density, and temperature conventions must be explicit.
\item Normalization affects numerical conditioning and machine-learning loss balance.
\end{itemize}

\SymbolsAcronymsSection
\begin{symtable}
$L_0,t_0,U_0$ & reference scales & Characteristic length, time, and velocity. \\
$V_A$ & Alfv'en speed & $B/\sqrt{\mu_0\rho}$. \\
$R_m$ & magnetic Reynolds number & Magnetic advection-to-diffusion ratio. \\
$M_A$ & Alfv'en Mach number & $U/V_A$. \\
$\hat Q$ & normalized variable & $Q/Q_0$. \\
\end{symtable}

\section{Exercises}
\begin{enumerate}[leftmargin=1.8em]
\item Derive the dimensions of $\epsilon_0$, electrical resistivity, current density, and electric potential.
\item Nondimensionalize $\partial_tu=D\partial_x^2u$ and derive the diffusion time.
\item Show that $B^2/(\mu_0\rho)$ has dimensions of velocity squared.
\item Convert $\omega=1.8\times10^6\ \si{s^{-1}}$ to hertz.
\item Construct a normalization table for the current Module 4 benchmark.
\item Explain how poor scaling can make a PINN satisfy one equation while violating another.
\end{enumerate}

\section{References}
\begin{thebibliography}{99}
\bibitem{ref1} G. B. Arfken, H. J. Weber, and F. E. Harris, \emph{Mathematical Methods for Physicists}, 7th ed., Academic Press (2013).
\bibitem{ref2} G. Strang, \emph{Linear Algebra and Its Applications}, 4th ed., Brooks/Cole (2006).
\bibitem{ref3} W. A. Strauss, \emph{Partial Differential Equations: An Introduction}, 2nd ed., Wiley (2007).
\bibitem{ref4} F. F. Chen, \emph{Introduction to Plasma Physics and Controlled Fusion}, 3rd ed., Springer (2016).
\bibitem{ref5} R. J. Goldston and P. H. Rutherford, \emph{Introduction to Plasma Physics}, CRC Press (1995).
\bibitem{ref6} P. M. Bellan, \emph{Fundamentals of Plasma Physics}, Cambridge University Press (2006).
\end{thebibliography}

\end{document}
